\unitheading{K1\quad 理想物理模型：质点}{unit-k01}
\addcontentsline{toc}{chapter}{K1 理想物理模型：质点}

自然界中的物体都处在运动过程中；所谓静止，只能相对于某个指定参考系而言。在经典力学的运动学描述中，不把“绝对静止”作为可直接判定的状态。学习力学的第一步，是建立一套能够准确描写运动的方法：先说明研究对象，再规定参考系和坐标，最后用位置随时间的变化表达运动。只有在运动学描述清楚之后，才有可能进一步研究运动状态变化的原因。

物理研究通常从最简单的对象开始。把一个具有确定质量、但形状和内部结构可以忽略的运动物体抽象成几何点，称为\textbf{质点}（mass point）。质点不是客观存在的“无大小物体”，而是针对具体问题建立的理想模型。

能否把物体看成质点，不能只比较物体自身大小，还要比较它的大小与研究尺度，并判断姿态、转动和形变是否影响所问物理量。例如：

\begin{itemize}
  \item 研究跳水运动员入水姿态时，身体各部分的位置关系直接影响结果，不能把运动员看成质点；研究他从北京到上海的行程时，身体尺寸相对于路程可以忽略，可用质点表示。
  \item 研究月球绕地球公转时，月球半径远小于地月距离，通常可把月球当成质点；研究探测器在月面着陆时，月球表面的曲率、地形和自转均可能重要，不能再用同一个简化。
  \item 研究乒乓球质心的抛体轨迹时可以暂时忽略球的大小；研究旋转引起的弧线或球与球拍接触时，必须保留半径和自转。
\end{itemize}

\EarthMoonScaleDiagram

建模时应明确写出“忽略了什么”。质点模型通常忽略物体的形状、姿态、转动和内部自由度，只保留质量与位置。若后续推导需要角速度、转动惯量、接触几何或形变能，就说明质点模型已经不足，必须升级为刚体、连续介质或可变形体模型。

判断模型是否合理的可靠方法，是做尺度和极限检查。设物体特征尺寸为 $\ell$、运动轨迹的特征尺度为 $L$。若 $\ell/L\ll1$，并且观测精度不要求分辨 $\ell$ 量级的差异，质点近似通常可靠；若结果对 $\ell$ 的一阶修正已经不可忽略，就不能仅凭“看起来很小”继续使用质点模型。
